\section{Passage choices in the opening example}
\label{app:example}
Table~\ref{tab:example-passages} gives the question in the introduction and the first sentence of each of the 10 passages. Rubin's and Sharapova's passages give the birth dates needed to determine which player is older.

\begin{table}[H]
\caption{The opening HotpotQA question, the instruction to choose every passage that supports the answer, and the 10 passage titles and texts in the order sent.}
\label{tab:example-passages}
\centering
\small
% Written by scripts/paper/paper_request.py sharapova_rubin --cut ". ", which ends each passage after its first
% sentence, and edited by hand: an ellipsis in place of the Russian name of Sharapova and its markup, and \& in place
% of &amp;. Running the script again overwrites the edits.
\begin{tabular}{@{}p{0.25\linewidth}p{0.68\linewidth}@{}}
\toprule
Question & who is older, Maria Sharapova or Chanda Rubin? \\
Instruction & Choose every passage in the list that supports the answer to the question. \\
\midrule
Title & Passage \\
\midrule
1997 IGA Classic -- Doubles & Chanda Rubin and Brenda Schultz-McCarthy were the defending champions but did not compete that year. \\
1997 EA-Generali Ladies Linz & The 1997 EA-Generali Ladies Linz was a women's tennis tournament played on indoor hard courts at the Intersport Arena in Linz in Austria that was part of Tier III of the 1997 WTA Tour.  \ldots \\
1997 Australian Open -- Women's Doubles & Chanda Rubin and Arantxa S\'{a}nchez Vicario were the defending champions, but competed this year with different partners.  \ldots \\
1997 State Farm Evert Cup -- Doubles & Chanda Rubin and Brenda Schultz-McCarthy were the defending champions but only Rubin competed that year with Mary Joe Fernandez. \\
Chanda Rubin & Chanda Rubin (born February 18, 1976) is an American tennis player.  \ldots \\
Maria Sharapova & Maria Yuryevna Sharapova (\ldots{} born April 19, 1987) is a Russian professional tennis player.  \ldots \\
1997 Bausch \& Lomb Championships -- Doubles & Chanda Rubin and Arantxa S\'{a}nchez Vicario were the defending champions but they competed with different partners that year, Rubin with Brenda Schultz-McCarthy and S\'{a}nchez Vicario with Larisa Savchenko. \\
2001 Bank of the West Classic -- Doubles & Chanda Rubin and Sandrine Testud were the defending champions, but none competed this year.  \ldots \\
1995 Schweppes Tasmanian International -- Doubles & Linda Harvey-Wild and Chanda Rubin were the defending champions but they competed with different partners that year, Harvey-Wild with Leila Meskhi and Rubin with Kristie Boogert. \\
2002 Wimbledon Championships -- Women's Doubles & Lisa Raymond and Rennae Stubbs were the defending champions, but lost in the quarterfinals to Anna Kournikova and Chanda Rubin. \\
\bottomrule
\end{tabular}

\end{table}

We send the same question and passages with three Choice instructions. They are ``Choose every passage in the list that supports the answer to the question'' (\emph{Supports}), ``Choose every passage in the list that is needed to answer the question'' (\emph{Needed}), and ``Choose every passage in the list that has information that answers part of the question'' (\emph{Part of the question}). We pair them with the yes/no questions ``Does this passage support the answer to the question?'', ``Is this passage needed to answer the question?'' and ``Does this passage have information that answers part of the question?''. Each yes/no question is followed by ``Passage:'' and the title and text of one passage.

\begin{table}[H]
\caption{Mean model probability. \emph{Distractor} is the highest probability among the other eight passages. Luna is GPT-6 Luna Decisions.}
\label{tab:example-wordings}
\centering
\resizebox{\linewidth}{!}{%
% Written by hand, not by a script. The table of tab:example-wordings in sections/appendix/example.tex. The % after
% \end{tabular} keeps the end of the file from adding a space inside the \resizebox.
\begin{tabular}{llrrr}
\toprule
Instruction & Model & Rubin & Sharapova & Distractor \\
\midrule
Supports & Jev & 0.93 & 0.07 & 0.00 \\
 & Liquid d1 & 0.94 & 0.06 & 0.00 \\
 & Luna & 0.93 & 0.02 & 0.01 \\
\midrule
Needed & Jev & 0.78 & 0.22 & 0.00 \\
 & Liquid d1 & 0.91 & 0.09 & 0.00 \\
 & Luna & 0.95 & 0.01 & 0.01 \\
\midrule
Part of the question & Jev & 0.77 & 0.23 & 0.00 \\
 & Liquid d1 & 0.95 & 0.05 & 0.00 \\
 & Luna & 0.87 & 0.01 & 0.03 \\
\bottomrule
\end{tabular}%
}
\end{table}

\begin{table}[H]
\caption{Mean probability of yes, with \emph{Distractor} as in Table~\ref{tab:example-wordings}.}
\label{tab:example-wordings-yes}
\centering
\resizebox{\linewidth}{!}{%
% Written by hand, not by a script. The table of tab:example-wordings-yes in sections/appendix/example.tex. The %
% after \end{tabular} keeps the end of the file from adding a space inside the \resizebox.
\begin{tabular}{llrrr}
\toprule
Instruction & Model & Rubin & Sharapova & Distractor \\
\midrule
Supports & Jev & 0.31 & 0.14 & 0.05 \\
 & Liquid d1 & 0.16 & 0.07 & 0.05 \\
 & Luna & 0.47 & 0.00 & 0.01 \\
\midrule
Needed & Jev & 0.61 & 0.40 & 0.05 \\
 & Liquid d1 & 0.53 & 0.41 & 0.03 \\
 & Luna & 1.00 & 0.87 & 0.80 \\
\midrule
Part of the question & Jev & 0.63 & 0.52 & 0.10 \\
 & Liquid d1 & 0.81 & 0.52 & 0.08 \\
 & Luna & 1.00 & 0.96 & 0.50 \\
\bottomrule
\end{tabular}%
}
\end{table}

Under each Choice instruction, every model assigns Rubin's passage a probability of at least 0.77. The probability assigned to Sharapova's passage is at most 0.23 (Table~\ref{tab:example-wordings}). The probability of yes for Sharapova's passage is 0.00 to 0.14 under Supports and 0.52 to 0.96 under Part of the question (Table~\ref{tab:example-wordings-yes}). The three instructions differ in meaning. This comparison does not separate the effect of wording from the effect of meaning.

Table~\ref{tab:example-invented-passages} gives the same list with invented names, tournaments, dates and locations. We send it with the Supports instruction and ask which player is older and which is younger. John is older than Jane. For the question about the older player, Jev and Liquid d1 favor John's passage and GPT-6 Luna Decisions favors Jane's. For the question about the younger player, all three models favor Jane's passage (Table~\ref{tab:example-invented}).

\begin{table}[H]
\caption{The questions and passages with invented names. John stands for Chanda Rubin and Jane for Maria Sharapova, and the passages are in the order of Table~\ref{tab:example-passages}.}
\label{tab:example-invented-passages}
\centering
\small
\begin{tabular}{@{}p{0.25\linewidth}p{0.68\linewidth}@{}}
\toprule
Question & who is older, Jane or John? \\
Question & who is younger, Jane or John? \\
Instruction & Choose every passage in the list that supports the answer to the question. \\
\midrule
Title & Passage \\
\midrule
1998 Harlow Classic – Doubles & John and Petra Lindqvist were the defending champions but did not compete that year. \\
1998 Varden Indoor Open & The 1998 Varden Indoor Open was a tennis tournament played on indoor hard courts at the Kestrel Arena in Varden in Switzerland that was part of Tier III of the 1998 Meridian Tour. The tournament was held from March 2 through March 8, 1998. John won the singles title. \\
1998 Corvale Open – Doubles & John and Ilse Marquart were the defending champions, but competed this year with different partners. John teamed up with Petra Lindqvist and were eliminated in third round, while Marquart teamed up with Dana Ferreira and lost in semifinals. \\
1998 Lakemont Cup – Doubles & John and Petra Lindqvist were the defending champions but only John competed that year with Nora Castellan. \\
John & John (born December 12, 1976) is an Italian tennis player. Winning eight Meridian Tour singles titles, he reached his highest ranking at World No. 5 on May 11, 1997, after reaching semifinals at the 1997 Corvale Open. John is also former World No. 7 in doubles, winning the Corvale Open in 1997 partnering with Ilse Marquart. \\
Jane & Jane (born January 13, 1990) is an Australian professional tennis player. A United States resident since 1996, Jane has competed on the Meridian Tour since 2005. She has been ranked world No. 1 in singles by the tour on four separate occasions, for a total of 17 weeks. She is one of nine players, and the only Australian, to hold the career Grand Slam. She is also a medalist of the Continental Games, having earned silver for Australia in singles at the 2014 Continental Games in Port Selwyn. \\
1998 Stonebridge Championships – Doubles & John and Ilse Marquart were the defending champions but they competed with different partners that year, John with Petra Lindqvist and Marquart with Vera Kolesnik. \\
2003 Bay Harbor Classic – Doubles & John and Camille Roussel were the defending champions, but none competed this year. John decided to focus on the singles tournament, while Roussel decided to rest after competing on the World Group Play-offs of the Nations Cup. \\
1996 Redcliff International – Doubles & Marta Hale-Winter and John were the defending champions but they competed with different partners that year, Hale-Winter with Sofia Nemec and John with Annika Brandt. \\
2004 Ashgrove Championships – Doubles & Laura Pemberton and Renee Stokes were the defending champions, but lost in the quarterfinals to Irina Voss and John. \\
\bottomrule
\end{tabular}

\end{table}

\begin{table}[H]
\caption{Model probability of John's and Jane's passages. Luna is GPT-6 Luna Decisions.}
\label{tab:example-invented}
\centering
\small
% Written by hand, not by a script. The table of tab:example-invented in sections/appendix/example.tex.
% The numbers are the answers of scripts/run/written_request.py to the requests john_jane and john_jane_younger on
% 2026-10-12, rounded to two decimals. The % after \end{tabular} keeps the end of the file from adding a space.
\begin{tabular}{llrrr}
\toprule
Question & Passage & Jev & Liquid d1 & Luna \\
\midrule
Who is older?   & John & 0.91 & 0.94 & 0.06 \\
                & Jane & 0.09 & 0.06 & 0.79 \\
Who is younger? & John & 0.02 & 0.15 & 0.03 \\
                & Jane & 0.98 & 0.85 & 0.89 \\
\bottomrule
\end{tabular}%

\end{table}
